\section{Bibliography}
% =============================================================================
% BIBLIOGRAPHY
% =============================================================================
% REQUIREMENT: Short bibliography identifying most relevant works
%
% STATUS: references.bib is comprehensive (30+ sources with comments)
%         This section can either list key works narratively OR be removed
%         since \bibliography{references} at end auto-generates the list
%
% DECISION NEEDED: Keep this section with narrative description of key sources,
%                  OR remove it and rely on auto-generated bibliography?
%
% IF KEEPING - structure as narrative of key sources by category:
%
% CATEGORY 1: IR THEORY (currently cited: waltz1979, wendt1992, barkin2010)
%   Key works in references.bib:
%   - Waltz 1979: Foundational neorealism
%   - Wendt 1992: Constructivist challenge
%   - Wendt 1999: Social Theory of IP (three cultures of anarchy)
%   - Mearsheimer 1994/2001: Offensive realism, critique of constructivism
%   - Barkin 2010: Synthesis attempt
%
% CATEGORY 2: ABM IN IR (not yet cited in text)
%   Key works in references.bib:
%   - Axelrod 1984/1997: Cooperation modeling
%   - Cederman 1997/2002: ABM for IR, balance of power
%   - Lustick 2000: ONLY constructivist ABM study (cite this!)
%   - Cederman & Daase 2003: "Constructivists need formal models"
%
% CATEGORY 3: RL/MARL (not yet cited in text)
%   Key works in references.bib:
%   - Leibo 2017: MARL in social dilemmas
%   - Dehkordi 2023: ML for agent specifications review
%   - Pan 2025: RL for IR interaction strategies
%
% CATEGORY 4: LLM AGENTS (not yet cited in text)
%   Key works in references.bib:
%   - Park 2023: Generative agents (landmark study)
%   - Hua 2023: WarAgent (LLMs simulating wars)
%   - Horiguchi 2024: Norm emergence in LLM agents
%   - Leng 2023: LLM social behavior, identity effects
%
% CATEGORY 5: METHODOLOGY (not yet cited in text)
%   Key works in references.bib:
%   - Larooij & Törnberg 2025: Generative ABM validation (CRITICAL)
%   - Unver 2018: Computational IR gap
%   - Pepinsky 2005: Critique of IR simulations
%
% TODO: Either write narrative paragraph for each category,
%       OR remove this section entirely (bibliography auto-generates)
% =============================================================================

\textbf{Core IR Theory}:
- Waltz -> Introduction Neorealism
- Wendt -> Introduction Constructivism, Criitism Neorealism
- Wendt -> Constructivism 'magnum opus'
- Mearsheimer -> Critism Constructivism, Revision Neorealism (Offensive Realism)

\textbf{ABM (in IR)}:
- Epstein
- Axelrod
- Cederman
- Lustick -> Only constructivist ABM?
- Cederman  Daase -> call for my thesis!!

\textbf{RL}:


\textbf{Social LLMs}
- Park
- Hua
- Horiguchi
- Leng
- Dai  Wu?

\textbf{Methodology}:
- Larooij Tornberg -> Generative ABM validation/


